\section*{\inTextCircled[10]{2}{15} Collapsed Tunnel}\stepcounter{section}\phantomsection\addcontentsline{toc}{section}{Collapsed Tunnel}
{\entryfont Farther along the forgotten road, the damage caused by the ancient collapse becomes increasingly severe. Great fractures split the surrounding masonry until the expedition eventually reaches a section where the ceiling has given way entirely, burying the route beneath an enormous mass of stone, timber, and the remains of old cargo.}
\begin{DndReadAloud}
	The passage ends abruptly beneath a mountain of debris. Broken blocks of masonry are wedged together with splintered beams, crushed carts, and centuries of accumulated dust, filling the tunnel from floor to ceiling.

	Here and there, narrow gaps disappear through the rubble toward darkness on the other side. A faint current of air passes between the stones, but no opening appears large enough for most of you to follow.

	Whatever road once continued beyond this point has been buried for centuries.
\end{DndReadAloud}
{\noindent\entryfont The collapse completely blocks the transport route. Although several narrow gaps extend through the rubble, none are large enough for most creatures to pass. Treat the mass of fallen stone and wreckage as the Ancient Debris described below.}
\begin{DndObject}[width=0.5\textwidth]{Ancient Debris}
	\DndObjectBasics[
		armor-class			= {10},
		hit-points			= {140},
		damage-threshold	= {10},
		damage-vulnerabilities	= {Fire, Force},
%		damage-resistances		= {Necrotic},
		damage-immunities		= {Necrotic, Piercing, Poison, Psychic, Radiant, Slashing},
	]
\end{DndObject}
{\noindent\entryfont Two old Kegs of Alchemist's Fire remain on the far side of the collapse. One rests immediately beside the debris, while the second stands farther along the tunnel and must first be brought closer. Each keg normally deals 6d6 Fire damage, but the debris takes maximum damage from an exploding keg. Detonating both beside the collapse is sufficient to clear the debris completely.

The characters must determine for themselves how to ignite the kegs.}

\subsection*{Clearing the Tunnel}\stepcounter{subsection}\phantomsection\addcontentsline{toc}{subsection}{Clearing the Tunnel}
{\entryfont The methods below describe several ways the expedition can clear the collapse using the characters' equipment and abilities. They are not exhaustive. Players should be encouraged to experiment with spells, tools, the surrounding environment, or their own solutions. Determine the resulting \textbf{Clearance Time} based on how quickly and carefully their approach would overcome the obstruction.}
\begingroup
	\DndSetThemeColor[PhbTan]
	\begin{DndComment}{DM Annotation - Clearance Time}
		Record whether the party achieves \textbf{Rapid}, \textbf{Delayed}, or \textbf{Prolonged} Clearance. The time spent overcoming the collapse affects the state of the ancient constructs encountered at the Control Point farther along the route.
	\end{DndComment}
\endgroup
\eject
\subsubsection*{Conjured Constructs}
\textbf{\textit{Rapid Clearance}}\\
{\entryfont L.O.R.E.-3 can expend a 3rd-level spell slot to cast Conjure Constructs. The summoned constructs tear through the unstable obstruction with overwhelming efficiency, clearing the route almost immediately. The process is highly destructive, however, and anything fragile buried within the collapse is crushed or lost.}
\subsubsection*{Drill and Alchemist's Fire}
\textbf{\textit{Delayed Clearance}}\\
{\entryfont B.R.E.A.C.H.-4 can use the Drill Arm to bore a narrow passage through the rubble. The opening is large enough for Small or smaller creatures to pass through, allowing B.R.E.A.C.H.-4 to reach the opposite side and move the second keg beside the collapse. If both kegs are subsequently detonated, the obstruction is destroyed.

The explosion leaves burning alchemical residue throughout the newly opened passage. The flames must be extinguished before the route can be crossed safely. Water, earth, appropriate magic, or any other reasonable method can put out the fire. A creature that moves through the burning passage without doing so takes 2d6 Fire damage.}
\subsubsection*{Crawling Through the Cracks}
\textbf{\textit{Delayed Clearance}}\\
{\entryfont Silas's Threadborn construction allows him to compress his frame enough to squeeze through the narrow openings between the fallen stones. Once on the other side, he can reach both kegs, though he lacks the strength to move the more distant one unaided.

The party must devise another method of bringing the second keg toward the collapse or otherwise clearing the obstruction. Ropes, magic, improvised mechanisms, or assistance found farther along the passage may all provide possible solutions.

If Silas continues farther into the tunnels before the obstruction is cleared, treat his arrival at the next location as Rapid Clearance regardless of how long the remainder of the expedition ultimately takes.}
\subsubsection*{Careful Tunnelling}
\textbf{\textit{Prolonged Clearance}}\\
{\entryfont B.R.E.A.C.H.-4 can instead use the Drill Arm for its intended purpose and gradually excavate a passage large enough for the entire expedition. This requires considerably more time but causes little additional damage to whatever has been buried within the collapse.}
\subsubsection*{Breaking Through}
\textbf{\textit{Prolonged Clearance}}\\
{\entryfont The expedition can simply attack the debris until it is destroyed. Bludgeoning attacks, powerful magic, and other suitably destructive effects can gradually clear a path, though the debris' Damage Threshold and immunities make some methods considerably less effective than others. Fragile objects buried within the rubble may be damaged or destroyed during the process.}

\subsection*{Buried Among the Rubble}\stepcounter{subsection}\phantomsection\addcontentsline{toc}{subsection}{Buried Among the Rubble}
{\entryfont The collapse contains the remains of cargo, tools, and personal belongings that were buried when the tunnel gave way centuries ago. How much survives depends largely on how the expedition clears the obstruction.
\paragraph*{Rapid Clearance} Nothing useful survives the process.
\paragraph*{Delayed Clearance} Roll once on the Random Finds table.
\paragraph*{Prolonged, Careful Excavation} Roll three times on the Random Finds table.
\paragraph*{Prolonged, Destructive Clearance} Roll once on the Random Finds table when the debris is reduced to half its Hit Points and once again when it is destroyed. Particularly fragile finds may be destroyed at the DM’s discretion.
\begin{DndTable}[header=Random Finds]{lX}
	d10	& Find	\\
	1	& A small pouch containing several worn Dwarven trade tokens, each stamped with a different workshop mark.\\
	2	& A cracked cargo seal bearing the faded emblem of an unknown merchant house.\\
	3	& A bundle of corroded but still serviceable precision tools wrapped in brittle cloth.\\
	4	& A delicate brass measuring instrument covered in markings that do not resemble the surrounding Dwarven script.\\
	5	& A finely made folding lens, remarkably clear despite its age and noticeably different in construction from the surrounding equipment.\\
	6	& A palm-sized articulated clockwork figurine. Its intricate workmanship appears unlike that of the ancient Dwarven automatons.\\
	7	& A reinforced container holding one intact flask of Alchemist's Fire.\\
	8	& A Greater Potion of Healing.\\
	9	& A Dwarven cargo tag bearing a second, much smaller inscription in an unfamiliar language.\\
	10	& A compact brass arcane component that still carries faint residual magic.\\
\end{DndTable}}